\documentclass[11pt]{article}
\usepackage[T1]{fontenc}
\usepackage{lmodern}
\usepackage[margin=1in]{geometry}
\usepackage{amsmath,amssymb,amsthm,mathtools,booktabs,array}
\usepackage{microtype,fancyhdr,enumitem}
\usepackage[hidelinks]{hyperref}
\hypersetup{pdftitle={IE-28: Rigorous Extensions and Exact Certificates},pdfauthor={ChatGPT},pdfsubject={Partial resolution; the arbitrary-node all-stage statement is not proved}}
\setlength{\headheight}{14pt}
\pagestyle{fancy}\fancyhf{}
\lhead{\small IE-28: rigorous extensions}\rhead{\small Scope-limited results}\cfoot{\thepage}
\setlength{\parindent}{0pt}\setlength{\parskip}{5pt}
\setlist{nosep,leftmargin=*}
\newtheorem{theorem}{Theorem}[section]
\newtheorem{lemma}[theorem]{Lemma}
\newtheorem{proposition}[theorem]{Proposition}
\newtheorem{corollary}[theorem]{Corollary}
\theoremstyle{remark}\newtheorem{remark}[theorem]{Remark}
\DeclareMathOperator{\diag}{diag}\DeclareMathOperator{\tr}{tr}
\newcommand{\R}{\mathbb R}\newcommand{\N}{\mathbb N}
\newcommand{\eps}{\varepsilon}
\newcommand{\coef}[2]{[#1^{#2}]}
\title{\textbf{IE-28: Rigorous Extensions\\and Exact Certificates}\\[6pt]
\large All-stage local existence, an explicit node family,\\and a limitation of the eigenpolynomial construction}
\author{Prepared by ChatGPT}
\date{12 September 2026}
\begin{document}
\maketitle\thispagestyle{empty}
\begin{abstract}
For the collocation integration matrix on distinct positive nodes, the target is
one positive diagonal $D$ for which $I-D^{-1}A$ is nilpotent. We give a
self-contained all-stage existence theorem for nodes sufficiently close to a
common positive point. The proof regularizes a determinant at coincident nodes,
constructs an algebraic confluent seed by a triangular recursion, and applies
the implicit function theorem with an explicitly nonsingular Jacobian.
We also characterize the node sets admitting a scalar diagonal: they are scaled
zeros of the ordinary Laguerre polynomial. Exact integer interval calculations
prove twelve additional point or uniform-neighborhood assertions in dimensions
four through seven. A further exact sign certificate shows that, at eleven
stages near confluence, no eigenpolynomial with all nonnegative coefficients
can produce the required diagonal, although a positive diagonal does exist.
The arbitrary-node, arbitrary-stage assertion is \emph{not proved or disproved}
here. The earlier complete two- and three-stage arguments are retained.
\end{abstract}

\begin{center}
\fbox{\begin{minipage}{0.91\textwidth}
\textbf{Status of the requested full solution.}
No full solution was obtained. In particular, the local all-stage theorem below
must not be read as a theorem for every admissible node set. This archive
contains additional rigorous results and checked computations, not a claimed
resolution of IE-28. No claim of priority or independent peer review is made.
\end{minipage}}
\end{center}

\section*{What is established}
\begin{tabular}{@{}p{0.26\textwidth}p{0.65\textwidth}@{}}
\toprule
Quantifiers & Result and evidence \\
\midrule
Every node pair or triple & Complete analytic proofs, Sections~\ref{sec:small} and the preserved prior writeup.\\
Every stage count; nodes near a common positive point & Analytic existence proof, Theorem~\ref{thm:local}. The neighborhood depends on the stage count.\\
Every stage count; scaled Laguerre nodes & Explicit scalar diagonal, Theorem~\ref{thm:laguerre}.\\
Six node sets and six node neighborhoods, $4\le n\le7$ & Exact integer interval certificates, Section~\ref{sec:cert}.\\
Eleven stages near confluence & Exclusion of a \emph{stronger ansatz}, not of IE-28, Theorem~\ref{thm:obstruction}.\\
Every admissible node set at every stage count & \textbf{Not established.}\\
\bottomrule
\end{tabular}

\bigskip
{\small\tableofcontents}
\newpage
\section{The exact target and useful reductions}\label{sec:target}
Let $n\ge2$ and $0<c_1<\cdots<c_n\le1$. Define
\begin{equation}\label{eq:A}
 \ell_j(t)=\prod_{k\ne j}\frac{t-c_k}{c_j-c_k},\qquad
 A_{ij}=\int_0^{c_i}\ell_j(t)\,dt.
\end{equation}
The target is
\begin{equation}\label{eq:target}
 D=\diag(d_1,\ldots,d_n),\quad d_i>0,\qquad
 (I-D^{-1}A)^n=0.
\end{equation}
It is the fixed-diagonal stiff-limit problem stated in IE-28~\cite{entry} and
in the conjecture of van der Houwen and de Swart~\cite[Section 3.2.1]{houwen}.
The MIN-SR-S treatment~\cite[Section 2.2.3]{sdc} does not supply a general
existence theorem. Nilpotence of $A-D$, a triangular preconditioner, or a
sequence of different diagonal sweep matrices would not solve
\eqref{eq:target}. No ordering of the $d_i$ is required.

Put $C=\diag(c_i)$ and $V_{ij}=c_i^{j-1}$. Interpolating monomials and integrating
exactly gives
\begin{equation}\label{eq:V}
 A=CV\diag(1,1/2,\ldots,1/n)V^{-1},\qquad
 \det A=\frac{\prod_i c_i}{n!}>0.
\end{equation}
All eigenvalues of $D^{-1}A$ are one if and only if all eigenvalues of its
inverse $A^{-1}D$ are one. Cayley--Hamilton then gives \eqref{eq:target}.
Equivalently, with $G=D^{-1}$,
\begin{equation}\label{eq:char}
 \det(I-\lambda AG)=(1-\lambda)^n.
\end{equation}
Here $AG$ and $GA$ are similar. Coefficient identities, rather than numerical
eigenvalues near a defective multiple eigenvalue, are used throughout.

For later computation let
\[
 w_i=\biggl(\prod_{j\ne i}(c_i-c_j)\biggr)^{-1},\qquad
 S=\diag(c_i/w_i).
\]
Direct differentiation of the interpolation formula gives
\begin{equation}\label{eq:M}
 A^{-1}=SMS^{-1},\qquad
 M_{ij}=\begin{cases}
 (c_i-c_j)^{-1},&i\ne j,\\
 c_i^{-1}+\displaystyle\sum_{j\ne i}(c_i-c_j)^{-1},&i=j.
 \end{cases}
\end{equation}
Indeed, if $W_{ij}=c_i^j$ and $W'_{ij}=j c_i^{j-1}$, then
$A^{-1}=W'W^{-1}$. Equivalently it maps values of a polynomial $q$ of degree
at most $n$, with $q(0)=0$, to the values of $q'$.
Applying this to the cardinal polynomials
$t\ell_j(t)/c_j$ proves \eqref{eq:M}. Diagonal matrices commute with $S$, so
$A^{-1}D$ is similar to $MD$.

Two immediate necessary identities are
\begin{equation}\label{eq:traceprod}
 \prod_i d_i=\frac{\prod_i c_i}{n!},\qquad
 \tr(MD)=\sum_i\frac{d_i}{c_i}
       +\sum_{i<j}\frac{d_i-d_j}{c_i-c_j}=n.
\end{equation}
They are not sufficient in general.

\begin{lemma}[Eigenpolynomial]\label{lem:eigen}
If $p$ is a nonzero real polynomial of degree at most $n$, $p(0)=0$, and
$p'(c_i)\ne0$, then $d_i=p(c_i)/p'(c_i)$ makes one an eigenvalue of
$A^{-1}D$, provided the vector of sampled values of $p$ is nonzero.
\end{lemma}
\begin{proof}
The polynomial-value identity is $A(p'(c_i))_i=(p(c_i))_i$.
Thus $D^{-1}A$ fixes the nonzero vector $(p'(c_i))_i$. The stated sampled
vector is nonzero whenever all $p(c_i)>0$, as in the applications below.
\end{proof}
Scaling causes no loss: $A(sc)=sA(c)$ for $s>0$, so $D\mapsto sD$ preserves
the target when the nodes are scaled by $s$.

\section{Complete small-stage cases}\label{sec:small}
This section summarizes the complete constructions from the preserved
\texttt{prior\_work/writeup.pdf}; their short proofs are included here.

\subsection{Two stages}
Let $c_1=a<c_2=b$ and $r=\sqrt{2ab}$. Set
\begin{equation}\label{eq:n2}
 d_i=\frac{c_i(c_i+r)}{2c_i+r},\qquad i=1,2.
\end{equation}
The polynomial $p(t)=t(t+r)$ gives eigenvalue one by Lemma~\ref{lem:eigen}.
Moreover
\[
 \frac{2d_1d_2}{ab}=1
 \quad\Longleftrightarrow\quad
 2(a+r)(b+r)=(2a+r)(2b+r)
 \quad\Longleftrightarrow\quad r^2=2ab.
\]
Thus $\det(A^{-1}D)=1$, so its other eigenvalue is also one. Positivity is
immediate, and Cayley--Hamilton proves the target.

\subsection{Three stages}
For $r>0$ and $0\le\theta\le1$, use
\begin{equation}\label{eq:n3ansatz}
 p(t)=t(t+r)(\theta t+r),\qquad
 d_i(r,\theta)=\left(\frac1{c_i}+\frac1{c_i+r}
                         +\frac{\theta}{\theta c_i+r}\right)^{-1}.
\end{equation}
There is always an eigenvalue one. Write
\[
 H_i(r,\theta)=1+\frac{c_i}{c_i+r}
                    +\frac{\theta c_i}{\theta c_i+r},\qquad
 G(r,\theta)=\det(MD)=6\prod_i H_i(r,\theta)^{-1}.
\]
For each fixed $\theta$, $G$ is strictly increasing in $r$ and tends to $6$ as
$r\to\infty$. Its limit as $r\downarrow0$ is $6/27$ for $\theta>0$, and
$6/8$ for $\theta=0$. Hence there is a unique $r(\theta)>0$ with $G=1$.

This root is continuous on $[0,1]$. To justify continuity at the corner
$(r,\theta)=(0,0)$, take any $r$ small enough that
$6\prod_i(1+c_i/(c_i+r))^{-1}<1$; then $G(r,\theta)<1$ for every $\theta$.
There is also a uniform upper bracket: at $r=4c_3$ every $H_i\le7/5$, so
$G\ge6/(7/5)^3>1$. Consequently the roots lie in a common compact subinterval
of $(0,4c_3)$; strict monotonicity and continuity give the claim.

Set $T(\theta)=\tr(MD(r(\theta),\theta))-3$ and let
$\sigma_1=c_1+c_2+c_3$, $\sigma_2=c_1c_2+c_1c_3+c_2c_3$,
$\sigma_3=c_1c_2c_3$. At $\theta=0$, substitution in
\eqref{eq:traceprod} and clearing denominators yields
\begin{equation}\label{eq:end0}
 T(0)=\frac{r(3r^2+3\sigma_1r+2\sigma_2)}{\prod_i(2c_i+r)}>0.
\end{equation}
At $\theta=1$, write $P_3=\prod_i(3c_i+r)$ and
\[
 F_3=6\prod_i(c_i+r)-P_3
    =5r^3+3\sigma_1r^2-3\sigma_2r-21\sigma_3.
\]
The determinant constraint is $F_3=0$, and direct simplification gives
\begin{equation}\label{eq:end1}
 P_3T(1)=F_3-2(r^3+3\sigma_3)<0.
\end{equation}
Both identities are verified as formal rational identities by the supplied
symbolic test. The intermediate value theorem supplies $\theta_*$ with $T=0$.
At this point $A^{-1}D$ has an eigenvalue one, trace three, and determinant
one. Its other two eigenvalues have sum two and product one, hence are both
one. This proves \eqref{eq:target} for every node triple.

Finally the denominator in \eqref{eq:n3ansatz} is strictly between $1/c_i$
and $3/c_i$, and it strictly decreases as $c_i$ increases. Therefore
\[
 c_i/3<d_i<c_i,\qquad d_1<d_2<d_3.
\]
The use of trace, determinant, and one known eigenvalue in this last spectral
step is special to three stages; it does not settle four or more stages.

\section{An explicit family at every stage count}\label{sec:laguerre}
\begin{lemma}[Characteristic polynomial from the node polynomial]\label{lem:chi}
Let $\pi(t)=\prod_{i=1}^n(t-c_i)$. Then
\begin{equation}\label{eq:chiA}
 \det(zI-A)=\frac1{n!}\sum_{j=0}^n\pi^{(j)}(0)z^j
 =\sum_{k=0}^n(-1)^k\frac{(n-k)!}{n!}e_k(c)z^{n-k},
\end{equation}
where $e_k$ is the $k$th elementary symmetric polynomial.
\end{lemma}
\begin{proof}
Identify sampled vectors with polynomials of degree at most $n-1$.
Integration followed by reduction modulo $\pi$ represents $A$. In the basis
$1,t,t^2/2!,\ldots,t^{n-1}/(n-1)!$, integration maps each basis vector other
than the last to the next one. For the last vector it gives $t^n/n!$, which,
modulo $\pi$, is
\[
 -\sum_{j=0}^{n-1}\frac{\pi^{(j)}(0)}{n!}\frac{t^j}{j!}.
\]
The representing matrix is a companion matrix with ones on the first
subdiagonal and these coefficients in its last column. Its characteristic
polynomial is the first expression in \eqref{eq:chiA}. The second follows by
expanding $\pi$.
\end{proof}

Define the ordinary Laguerre polynomial by
\[
 L_n(x)=\sum_{j=0}^n(-1)^j\binom nj\frac{x^j}{j!}.
\]
\begin{theorem}[Scalar diagonals are exactly the scaled Laguerre case]\label{thm:laguerre}
Let $x_1<\cdots<x_n$ be the zeros of $L_n$. They are positive and simple.
For any $\delta>0$ with $\delta x_n\le1$, the nodes $c_i=\delta x_i$ satisfy
\eqref{eq:target} with
\begin{equation}\label{eq:scalar}
 D=\delta I.
\end{equation}
Conversely, if a node set admits a scalar positive diagonal $D=\delta I$
satisfying \eqref{eq:target}, its node polynomial is
$\pi(t)=(-\delta)^n n!L_n(t/\delta)$.
\end{theorem}
\begin{proof}
For completeness, the Rodrigues identity
\[
 e^{-x}L_n(x)=\frac1{n!}\frac{d^n}{dx^n}(e^{-x}x^n)
\]
follows by the product rule. Integrating by parts $n$ times shows that $L_n$
is orthogonal to every polynomial of degree less than $n$ against $e^{-x}$
on $(0,\infty)$. The boundary terms vanish. If $L_n$ had fewer than $n$
sign-changing zeros there, multiply the factors corresponding to those zeros
to obtain a polynomial of degree less than $n$ whose product with $L_n$ has
one sign and is not zero almost everywhere. Its integral cannot vanish,
contradicting orthogonality. Thus all $n$ zeros are positive and simple.

For the scaled zeros, the monic node polynomial is
$\pi(t)=(-\delta)^n n!L_n(t/\delta)$, so
\[
 \frac{\pi^{(j)}(0)}{n!}=(-1)^{n+j}\delta^{n-j}\binom nj.
\]
Lemma~\ref{lem:chi} gives $\det(zI-A)=(z-\delta)^n$.
Cayley--Hamilton proves \eqref{eq:scalar}. Conversely the scalar target forces
that characteristic polynomial, so \eqref{eq:chiA} determines every Taylor
coefficient of $\pi$ and gives the stated Laguerre expression.
\end{proof}
This is an exact all-stage class, not a reduction of arbitrary nodes to this
class: scaling alone cannot turn an arbitrary node set into Laguerre zeros.

\section{Removing the singularity at coincident nodes}\label{sec:regularization}
For a polynomial $g$ of degree at most $n-1$ define, initially at distinct nodes,
\begin{equation}\label{eq:Phi}
 \Phi(\lambda;c,g)=
 \frac{\det[jc_i^{j-1}-\lambda g(c_i)c_i^j]_{i,j=1}^n}
 {n!\,\Delta(c)},\qquad
 \Delta(c)=\prod_{i<j}(c_j-c_i).
\end{equation}
The numerator is alternating in the node variables. In the polynomial ring
whose other indeterminates are $\lambda$ and the coefficients of $g$, it is
divisible by each $c_j-c_i$, hence by their product. Consequently $\Phi$ is a
polynomial and extends across all node collisions.

\begin{lemma}\label{lem:Phi}
At distinct positive nodes, if $G=\diag(g(c_i))$, then
\begin{equation}\label{eq:Phieq}
 \Phi(\lambda;c,g)=\det(I-\lambda AG).
\end{equation}
Thus $g(c_i)>0$ and $\Phi=(1-\lambda)^n$ give a solution
$D=\diag(1/g(c_i))$ of IE-28.
\end{lemma}
\begin{proof}
Use $W,W'$ from Section~\ref{sec:target}. The numerator is
$\det(W'-\lambda GW)$. Since $A^{-1}=W'W^{-1}$ and
$\det W'=n!\Delta(c)$, its ratio to $\det W'$ is
$\det(A^{-1}-\lambda G)/\det A^{-1}$, which is \eqref{eq:Phieq}.
\end{proof}

A division-free formula is useful both as a check and as a numerical
implementation. Let $h_m(x_1,\ldots,x_r)$ denote the complete homogeneous
symmetric polynomial, with $h_0=1$ and $h_m=0$ for $m<0$. The elementary
monomial divided-difference identity is
\[
 [x_1,\ldots,x_{k+1}]t^m=h_{m-k}(x_1,\ldots,x_{k+1}).
\]
For $g(t)=\sum_{r=0}^{n-1}g_rt^r$, define, for $0\le k\le n-1$,
$1\le j\le n$,
\begin{align}\label{eq:pencil}
 (B_0)_{k,j-1}&=j h_{j-1-k}(c_1,\ldots,c_{k+1}),\\
 (B_1)_{k,j-1}&=\sum_{r=0}^{n-1}g_rh_{r+j-k}(c_1,\ldots,c_{k+1}).\nonumber
\end{align}
Changing evaluation rows into successive divided-difference rows divides their
determinant by $\Delta(c)$. Hence
\begin{equation}\label{eq:pencilidentity}
 \Phi(\lambda;c,g)=\det(B_0-\lambda B_1)/n!,\qquad \det B_0=n!.
\end{equation}
Unlike the defining quotient, this formula can be evaluated at coincident
nodes without subtracting nearly equal nodes or dividing by node gaps.

\section{All-stage existence near a common positive point}\label{sec:local}
Parameterize $g$ by $a=(a_1,\ldots,a_n)$ through
\begin{equation}\label{eq:Ha}
 H_a(z)=\sum_{j=1}^n a_jz^j,\qquad
 g_a(t)=-\sum_{j=1}^n j a_j(1-t)^{j-1}.
\end{equation}
In particular $g_a(1)=-a_1$ and $\int_1^t g_a(s)\,ds=H_a(1-t)$.

\begin{lemma}[Confluent determinant]\label{lem:confluent}
At $c_1=\cdots=c_n=1$, the polynomial extension satisfies
\begin{align}\label{eq:conf}
 \Phi(\lambda;\mathbf1,g_a)
 &=\coef{z}{n}\frac{\exp(\lambda H_a(z))}{1-z}
 =1+\sum_{k=1}^n f_k(a)\lambda^k,\\
 f_k(a)&=\frac1{k!}\coef{z}{n}\frac{H_a(z)^k}{1-z}.\nonumber
\end{align}
Only finite Taylor coefficients are used in this expression.
\end{lemma}
\begin{proof}
The confluent evaluation determinant in \eqref{eq:Phi} has rows consisting
of the Taylor coefficients of $q'-\lambda g_aq$ of orders $0,\ldots,n-1$
at $t=1$, with $q$ ranging over $t,t^2,\ldots,t^n$.
Write $s=t-1$ and $q(t)=\sum_{m=0}^n y_ms^m$. The constraint $q(0)=0$ is
\begin{equation}\label{eq:boundary}
 \sum_{m=0}^n(-1)^m y_m=0.
\end{equation}
For an arbitrary $y_0$, the equations that the indicated residual coefficients
vanish form a triangular recursion for $y_1,\ldots,y_n$, with pivots
$1,2,\ldots,n$. Its solution is
\[
 y_m=y_0\coef{s}{m}\exp(\lambda H_a(-s)),\qquad 0\le m\le n.
\]
Eliminating $y_1,\ldots,y_n$ from these equations and the boundary equation
leaves $y_0$ multiplied by
\[
 \sum_{m=0}^n\coef{z}{m}\exp(\lambda H_a(z)).
\]
To see that this is an identity of determinants, not just an equivalence of
zero sets, order the augmented equations with \eqref{eq:boundary} first and
the residual equations next. Their determinant is $n!$ times the displayed
factor by the triangular elimination. Eliminating $y_0$ using its unit
coefficient in \eqref{eq:boundary} gives precisely the restricted residual
matrix. The change from coefficients in $t,t^2,\ldots,t^n$ to
$y_1,\ldots,y_n$ has determinant one (its entries are $\binom jk$).
Division by $n!$ therefore gives the normalized determinant in
\eqref{eq:Phi}. Summing the Taylor coefficients is coefficient extraction
against $(1-z)^{-1}$, proving \eqref{eq:conf}.
\end{proof}

\begin{lemma}[Triangular algebraic seed]\label{lem:seed}
There is a unique $a^*\in\R^n$ with $a_1^*<0$ such that
$f_k(a^*)=(-1)^k\binom nk$ for $k=1,\ldots,n$. It is given by
\begin{align}\label{eq:seed}
 r_n&=(n!)^{1/n},\qquad a_1^*=-r_n,\\
 a_j^*&=\frac{(k-1)!}{(a_1^*)^{k-1}}
 \left((-1)^k\binom nk-f_k(a_1^*,\ldots,a_{j-1}^*,0,\ldots,0)\right),
 \quad k=n-j+1,\quad 2\le j\le n.\nonumber
\end{align}
The Jacobian of $a\mapsto(f_n,f_{n-1},\ldots,f_1)$ at this seed is
nonsingular, with determinant
\begin{equation}\label{eq:Jac}
 \frac{(a_1^*)^{n(n-1)/2}}{\prod_{m=0}^{n-1}m!}\ne0.
\end{equation}
\end{lemma}
\begin{proof}
Since $H_a$ has no constant term, $f_n=a_1^n/n!$, and its required value fixes
the unique negative $a_1=-r_n$. For $k=n-j+1$, the terms of $H_a^k$ of degree
at most $n$ involve no coefficient beyond $a_j$. A term containing $a_j$
can have degree at most $n$ only when it contains $a_jz^j$ exactly once and
$a_1z$ in all other factors. For $j\ge2$ its coefficient in $f_k$ is therefore
$a_1^{k-1}/(k-1)!$, nonzero when $a_1\ne0$. This proves the recursion.
In the reversed row order, the Jacobian is triangular, and its diagonal is
\[
 \frac{a_1^{n-1}}{(n-1)!},\frac{a_1^{n-2}}{(n-2)!},\ldots,a_1,1.
\]
Their product proves \eqref{eq:Jac}. The first diagonal entry comes by
differentiating $f_n$ and is consistent with this expression.
\end{proof}

\begin{theorem}[All-stage clustered-node existence]\label{thm:local}
For every integer $n\ge2$, there exists $\eps_n>0$ such that all distinct
nodes with $1-\eps_n<c_1<\cdots<c_n\le1$ admit a positive diagonal satisfying
\eqref{eq:target}. There is a real-analytic branch, in polynomial coordinates,
with
\begin{equation}\label{eq:limit}
 d_i\longrightarrow(n!)^{-1/n}\qquad(c\longrightarrow\mathbf1).
\end{equation}
More generally, nodes sufficiently close to any common $c_*>0$ admit such a
diagonal, with limiting entry $c_*(n!)^{-1/n}$.
\end{theorem}
\begin{proof}
Let $\mathcal F(a,c)$ be the vector of coefficients of $\lambda^n,\ldots,
\lambda$ in $\Phi(\lambda;c,g_a)-(1-\lambda)^n$.
By Section~\ref{sec:regularization}, this is a polynomial map in $(a,c)$,
including at $c=\mathbf1$. Lemmas~\ref{lem:confluent} and~\ref{lem:seed} give
$\mathcal F(a^*,\mathbf1)=0$ and an invertible derivative with respect to $a$.
The real-analytic implicit function theorem yields a solution $a(c)$ for all
$c$ in some neighborhood of $\mathbf1$.

At the center, $g_{a^*}(1)=r_n>0$. By continuity, after shrinking that
neighborhood, $g_{a(c)}(c_i)>0$ for every $i$. Set $d_i=1/g_{a(c)}(c_i)$.
Lemma~\ref{lem:Phi} proves \eqref{eq:target} at all distinct positive nodes
in the neighborhood. It also gives the limit \eqref{eq:limit}. Scaling from
Section~\ref{sec:target} proves the assertion at a common point $c_*$.
\end{proof}

\begin{remark}[What locality means]
The theorem gives a positive neighborhood for each fixed $n$. It does not
bound $\eps_n$ uniformly in $n$, does not claim that $\eps_n=1$, and does not
prove that this branch extends through the entire ordered-node domain.
The coincident-node matrix $A$ itself is not used: the implicit function
theorem is applied to its polynomial determinant regularization.
\end{remark}

\section{An exact obstruction to a stronger ansatz}\label{sec:obstruction}
A tempting route is to choose
\begin{equation}\label{eq:strong}
 p(t)=t\prod_{j=1}^{n-1}(t+\alpha_j),\qquad\alpha_j>0,
 \quad d_i=p(c_i)/p'(c_i).
\end{equation}
It ensures positivity and $c_i/n<d_i<c_i$, but those extra restrictions were
not required by IE-28. The following result prevents promoting this route to
an unrestricted all-stage proof.

Define the confluent normalized polynomial
\begin{equation}\label{eq:pstar}
 p_*(t)=\sum_{m=0}^n\left(\coef{z}{m}e^{H_{a^*}(z)}\right)(1-t)^m.
\end{equation}
Then $p_*(1)=1$, and $p_*(0)=\Phi(1;\mathbf1,g_{a^*})=0$.
It is the unique polynomial of degree at most $n$ with value one at $1$
whose residual $p_*'-g_{a^*}p_*$ has zero Taylor coefficients of orders
$0,\ldots,n-1$ there.

\begin{proposition}[Exact integer sign certificate]\label{prop:sign}
For $n=11$, the coefficient $\beta_{11}=\coef{t}{11}p_*(t)$ satisfies
\begin{equation}\label{eq:negative}
 -1.664575\,10^{-10}<\beta_{11}<-1.664574\,10^{-10}<0.
\end{equation}
These inequalities are verified with integer arithmetic, not inferred from
floating-point roots.
\end{proposition}
\begin{proof}[Certificate and its mathematical justification]
The file \texttt{cluster/certify\_seed.py} brackets $r=(11!)^{1/11}$ between
consecutive multiples of $10^{-100}$ by comparing integer eleventh powers.
The strict bracket is checked exactly. It propagates this bracket through
\eqref{eq:seed} and the truncated exponential in \eqref{eq:pstar}, using
outward-rounded dyadic intervals at 512 bits. All sums, products, reciprocals,
and factorial divisions enclose the exact real operations. The denominators
in \eqref{eq:seed} are nonzero powers of the strictly negative $a_1$.
Consequently the resulting interval encloses the exact algebraic coefficient
$(-1)^{11}\coef{z}{11}e^{H_{a^*}(z)}$. Its two endpoints are checked, by
integer cross multiplication, to lie strictly between the rational endpoints
in \eqref{eq:negative}. The full rational endpoints and the root bracket are
saved in \texttt{cluster/seed11\_certificate.json}. This finite verification
uses Python's standard library only. It is a reproducible computer-assisted
sign proof; the decimal display printed by the program is not used by the
proof.
\end{proof}

\begin{theorem}[No nonnegative-coefficient eigenpolynomial near confluence]\label{thm:obstruction}
For $n=11$ there exists $\eta>0$ such that, for any distinct nodes
$1-\eta<c_1<\cdots<c_{11}\le1$, no polynomial $p$ with
\[
 p(0)=0,\quad p\ne0,\quad \deg p\le11,
 \quad\text{all coefficients of $p$ nonnegative}
\]
can yield a solution of \eqref{eq:target} by $d_i=p(c_i)/p'(c_i)$.
In particular, \eqref{eq:strong} cannot solve these node sets. Nevertheless,
for a sufficiently small such $\eta$, positive diagonals \emph{do} exist by
Theorem~\ref{thm:local}.
\end{theorem}
\begin{proof}
Suppose to the contrary that there are node vectors $c^{(m)}\to\mathbf1$ and
such polynomials $p_m$ producing solutions. Normalize $p_m(1)=1$, which does
not change $p_m/p_m'$. Their coefficients form a compact simplex. Along a
subsequence they converge coefficientwise to a polynomial $p$ with
nonnegative coefficients, $p(0)=0$, and $p(1)=1$.

On a fixed real neighborhood of $1$, the $p_m$ are uniformly bounded away
from zero; for example $p_m(t)\ge2^{-11}$ for $1/2\le t\le1$.
Their derivatives of all fixed orders converge, so $g_m=p_m'/p_m$ and their
fixed-order derivatives converge uniformly on a smaller neighborhood.
Let $\widetilde g_m$ be the polynomial of degree at most ten interpolating
$g_m$ at $c^{(m)}$. Divided differences converge to derivatives at a repeated
node; this follows directly from the integral form of a divided difference
and the uniform derivative convergence. Thus $\widetilde g_m$ converges,
coefficientwise, to the degree-ten Taylor polynomial of $g=p'/p$ at $1$.
Call this limit $g_0$.

The solution assumption and Lemma~\ref{lem:Phi} give
$\Phi(\lambda;c^{(m)},\widetilde g_m)=(1-\lambda)^{11}$.
Polynomial continuity gives the same identity at $(\mathbf1,g_0)$.
Moreover $g_0(1)=p'(1)>0$; in fact $1\le p'(1)\le11$ from the normalized
nonnegative coefficients. By the uniqueness in Lemma~\ref{lem:seed},
$g_0=g_{a^*}$. The residual $p'-g_0p$ therefore vanishes to order at least
11 at $1$. Its Taylor recursion, together with $p(1)=1$, forces
$p=p_*$ from \eqref{eq:pstar}. This contradicts the negative leading
coefficient in Proposition~\ref{prop:sign}.
\end{proof}

The theorem does not say that a numerical solver failing to find negative
roots has disproved IE-28. It proves failure of a specified auxiliary
construction and separates that failure from the positive-diagonal question.

\section{Exact certificates away from confluence}\label{sec:cert}
The analytical local theorem does not identify its neighborhood size.
Independently, the archive supplies exact certificates for specified node
sets and explicit node boxes. No assertion in this section is based only on
a small residual.

For $M(c)$ from \eqref{eq:M}, write
\begin{equation}\label{eq:F}
 F_k(d,c)=\sum_{\substack{S\subseteq\{1,\ldots,n\}\\|S|=k}}
              \det M(c)[S,S]\prod_{i\in S}d_i-\binom nk,
 \quad 1\le k\le n.
\end{equation}
Then $F=0$ is exactly the characteristic-coefficient system for $MD$ having
all eigenvalues one. The Jacobian used by the verifier is the derivative of
this displayed polynomial with respect to $d$.

\begin{theorem}[Uniform contraction certificate]\label{thm:certificate}
Let $\mathcal C$ be a closed node box consisting only of separated positive
nodes with last node at most one, and let
$X=\{d:\|d-d^0\|_\infty\le r\}$ have strictly positive coordinates.
Let $R$ be a fixed real $n\times n$ matrix. Suppose rigorous bounds satisfy
\begin{align}\label{eq:bounds}
 \|RF(d^0,c)\|_\infty&\le\xi &&(c\in\mathcal C),\\
 \|I-RF_d(d,c)\|_\infty&\le q<1 &&(d\in X,\ c\in\mathcal C),\nonumber\\
 \xi+qr&<r.&&\nonumber
\end{align}
For every $c\in\mathcal C$ there exists exactly one zero of $F(\cdot,c)$ in
$X$. It gives a positive diagonal solving \eqref{eq:target}.
\end{theorem}
\begin{proof}
For each fixed $c$, define $T_c(d)=d-RF(d,c)$. The derivative bound and the
mean-value integral on the convex box $X$ make $T_c$ a contraction with
constant at most $q$. The first and third bounds show $T_c(X)\subset X$.
Banach's fixed-point theorem gives a unique fixed point in $X$.
For any one $d\in X$, $\|I-RF_d(d,c)\|<1$ implies that $RF_d(d,c)$ is
invertible, so the square matrix $R$ is invertible. Fixed points of $T_c$
are therefore exactly zeros of $F(\cdot,c)$. Positivity is part of the
hypothesis on $X$, and the coefficient characterization proves the target.
Uniqueness is only asserted \emph{inside this diagonal box}.
\end{proof}

\subsection{What the verifier checks}
All centers, radii, and preconditioner entries in the JSON certificates are
exact rational numbers (decimal strings denote exact terminating rationals).
The verifier uses intervals $[L/2^b,U/2^b]$ with integer endpoints. Addition,
multiplication, and reciprocals are rounded outward using integer floor and
ceiling division. Principal determinants are enclosed by division-free
subset dynamic programming, not by a floating-point factorization.
At $b=320$, it obtains integer upper bounds $E,Q$ for $2^b\xi,2^bq$.
For $r=r_N/r_D$, its final self-mapping test is precisely
\[
 E r_D+Q r_N<2^b r_N,
\]
and its contraction test is $Q<2^b$. Both are integer comparisons. The
positivity of the diagonal box, separation of every node interval, and the
upper node bound are also checked. A passing result is an application of
Theorem~\ref{thm:certificate}; the approximate candidate generator is outside
the trusted verification path.

\subsection{Certified node neighborhoods}
In the following table the last node is fixed at one; every other node may
vary independently by $\pm\delta$ about its listed center. The diagonal
box is centered at the rational vector in the certificate file and has
radius $r$. The displayed $q,\xi$ bounds have been rounded \emph{upward}
for readability; the actual acceptance tests use their exact rational
versions.

\begin{center}\small
\begin{tabular}{@{}clrrrr@{}}
\toprule
$n$ & Node center & $\delta$ & $r$ & $q$ upper & $\xi$ upper\\
\midrule
4 & $(.1,.3,.6,1)$ & $10^{-5}$ & $10^{-3}$ & $.222$ & $.000600$\\
4 & $c_i=i/4$ & $10^{-4}$ & $10^{-2}$ & $.475$ & $.00436$\\
5 & $c_i=i/5$ & $10^{-6}$ & $10^{-3}$ & $.139$ & $.000108$\\
6 & $c_i=i/6$ & $10^{-6}$ & $10^{-3}$ & $.346$ & $.000252$\\
7 & $c_i=i/7$ & $10^{-7}$ & $10^{-3}$ & $.771$ & $.0000622$\\
5 & $(.01,.02,.03,.04,1)$ & $10^{-15}$ & $10^{-8}$ & $.228$ & $9.0\,10^{-10}$\\
\bottomrule
\end{tabular}
\end{center}
All six boxes pass. There is also a separate, much tighter, point certificate
at each of these centers, using $r=10^{-30}$. Thus twelve certificates pass
in total. The separated five-node example needs a small box under the
straightforward interval bounds used here; no claim of a maximal neighborhood
is made. Scaling the nodes and diagonal extends each normalized assertion to
its corresponding positive scaling family.

\section{Continuation, checks, and the remaining global gap}\label{sec:gap}
\subsection{Numerical continuation is explicitly not a proof of global existence}
The regularization supplies a natural numerical path
\[
 c_i(s)=1-s(1-c_i),\quad0\le s\le1,
\]
starting at the exact algebraic seed. The supplied program performs damped
Newton corrections while keeping all sampled $g_a(c_i(s))$ positive.
For improved conditioning it uses $z_i=1-c_i$ and the polynomial basis
$q_j(t)=(1-t)^j-1$, which automatically satisfies the boundary condition.
In this basis the divided-difference derivative matrix is
\[
 T_{k,j-1}=-j h_{j-1-k}(z_1,\ldots,z_{k+1}),
\]
and the derivative of the $g_aq_j$ matrix with respect to $a_{r+1}$ is
\[
 -(r+1)\left(h_{r+j-k}(z_1,\ldots,z_{k+1})
                     -h_{r-k}(z_1,\ldots,z_{k+1})\right).
\]
The normalized pencil is $I-\lambda K(a)$ after left multiplication by
$T^{-1}$. Its coefficients and their derivatives are evaluated through
finite matrix recurrences, not by differentiating numerical eigenvalues.

The saved equally spaced examples for $n=8,11,15$ reach $s=1$ with positive
sampled $g$, using 60, 80, and 90 decimal working digits respectively.
Their final coefficient residuals are approximately $1.6\,10^{-38}$,
$4.3\,10^{-49}$, and $7.7\,10^{-49}$. These are high-precision
\emph{numerical examples only}; they are not among the twelve interval
certificates. A flag named \texttt{pass} in those numerical JSON files means
that the numerical stopping tests passed, not that a theorem was proved.

\subsection{Independent checks included}
The exact regression suite performs 1,500 rational interval-arithmetic trials,
50 rational determinant comparisons, formal seed-Jacobian checks through
$n=6$, confluent determinant checks through $n=5$, exact collocation
characteristic checks through $n=7$, and symbolic Laguerre identities through
$n=15$. It reruns all twelve certificates and confirms rejection of three
intentionally corrupted inputs. The previous two- and three-stage symbolic
identities are checked separately. These tests help detect implementation
errors; finite tests are not substitutes for the analytical all-stage proofs.

\subsection{What would still be needed for a full solution}
The missing implication is global:
\[
 \forall n\ge4\ \forall\,0<c_1<\cdots<c_n\le1,
 \quad\exists d\in(0,\infty)^n:\ F(d,c)=0.
\]
For the particular continuation route above, one would need a theorem that
prevents loss of the required solution before $s=1$: for example, a valid
continuation or degree argument controlling finite escape, approach to the
boundary of positivity, and singularities of the coefficient Jacobian.
Those conditions have not been established here. A nonsingular Jacobian at
$s=0$, or successful Newton tracking for finitely many target node sets,
does not establish them.

Several shortcuts therefore remain invalid. The trace and determinant
constraints are insufficient beyond the small-stage argument; generic
complex multiplicative inverse-eigenvalue theorems~\cite{rosenthal} do not
guarantee positive real diagonals; and the nonnegative-coefficient
construction is ruled out near an eleven-stage confluence by
Theorem~\ref{thm:obstruction}. Earlier low-order algebraic and higher-order
numerical investigations~\cite{lioen} are valuable context, not an all-node
proof. The bounded literature search conducted for this attempt did not
locate a full theorem or a counterexample settling the stated target.

\section{Reproduction and archive map}
The main proofs are in this document and its editable LaTeX source.
\texttt{README.md} gives commands for each check. The key commands, run from
the extracted archive root, are:
\begin{verbatim}
python certification/verify.py \
  certification/point_certificates.json \
  certification/neighborhood_certificates.json \
  --report results/verification_report.json
python cluster/certify_seed.py \
  --report cluster/seed11_certificate.json
python tests/test_exact.py
python prior_work/scripts/check_symbolic.py
\end{verbatim}
The first two use the Python standard library only. The symbolic regression
requires SymPy. Numerical continuation and candidate generation require
mpmath, and the preserved previous numerical examples also use NumPy.
Dependencies and the tested environment are recorded in the archive.
\texttt{prior\_work/} preserves the preceding partial writeup and its sources;
its status statements describe that earlier version, not the additional
results proved in the present document.

\begin{thebibliography}{9}
\bibitem{entry}
\emph{OpenProblemsInNLA}, IE-28, ``Positive diagonal preconditioning with a
nilpotent stiff limit.'' Public problem statement, accessed 12 September 2026;
the entry records a bounded status check dated 11 September 2026.
\url{https://github.com/ajt60gaibb/OpenProblemsInNLA/blob/main/linear-systems-and-elimination/IE-28/README.md}.

\bibitem{houwen}
P.~J. van der Houwen and J.~J.~B. de Swart,
``Triangularly implicit iteration methods for ODE-IVP solvers,''
\emph{SIAM Journal on Scientific Computing} 18(1) (1997), 41--55.
DOI: \href{https://doi.org/10.1137/S1064827595287456}{10.1137/S1064827595287456}.
Open copy: \url{https://ir.cwi.nl/pub/2191/2191D.pdf}.

\bibitem{sdc}
G.~\v{C}aklovi\'{c}, T.~Lunet, S.~G\"otschel, and D.~Ruprecht,
``Improving efficiency of parallel across the method spectral deferred
corrections,'' \emph{SIAM Journal on Scientific Computing} 47(1) (2025),
A430--A453. DOI: \href{https://doi.org/10.1137/24M1649800}{10.1137/24M1649800}.
Open text: \url{https://epubs.siam.org/doi/10.1137/24M1649800}.

\bibitem{lioen}
W.~M. Lioen, ``On the diagonal approximation of full matrices,'' CWI report
NM-R9518 (1995); journal version in \emph{Journal of Computational and Applied
Mathematics} 75(1) (1996), 35--42.
Report record: \url{https://ir.cwi.nl/pub/4925};
PDF: \url{https://ir.cwi.nl/pub/4925/04925D.pdf}.

\bibitem{rosenthal}
J.~Rosenthal and X.~Wang, ``The multiplicative inverse eigenvalue problem over
an algebraically closed field,'' \emph{SIAM Journal on Matrix Analysis and
Applications} 23(2) (2001), 517--523.
Preprint: \url{https://arxiv.org/abs/math/0009163}.
Its algebraically closed-field setting is essential and is not a
positive-real existence theorem for IE-28.
\end{thebibliography}
\end{document}
